% TOP_PROBLEM: {"id": 2931, "title": "Kirby Problem 4.55", "category_id": 7, "category": {"id": 7, "name": "topology", "display_name": "Topology", "description": "Properties preserved under continuous deformations.", "slug": "topology", "order_index": 7, "created_at": "2026-07-31T15:26:25.671Z"}, "status": "open", "problem_number": "KP-4.55", "set_id": 11}
% FIRST_POSTED: 2026-10-01
% ENTRY_KIND: statement_only
% UnsolvedMath ID 2931; source catalogue code KP-4.55
% !TeX program = lualatex
% Definition and sources only; this file is not an LLM solution attempt.
\documentclass[11pt,a4paper]{article}
\usepackage[margin=25mm]{geometry}
\usepackage{fontspec}
\setmainfont{lmroman10-regular.otf}[BoldFont=lmroman10-bold.otf,
 ItalicFont=lmroman10-italic.otf,BoldItalicFont=lmroman10-bolditalic.otf]
\usepackage{amsmath,amssymb,mathtools}
\usepackage{unicode-math}
\setmathfont{latinmodern-math.otf}
\AtBeginDocument{\let\setminus\smallsetminus}
\usepackage{xurl,hyperref}
\hypersetup{colorlinks=true,urlcolor=blue}
\setcounter{secnumdepth}{0}
\setlength{\parindent}{0pt}
\setlength{\parskip}{5pt}
\setlength{\emergencystretch}{3em}
\begin{document}
\section*{Kirby Problem 4.55}\label{2931}
{\small UnsolvedMath ID 2931 · KP-4.55 · Topology · Kirby's Problems in Low-Dimensional Topology}
\subsection{Definitions and mathematical statement}
For a closed connected oriented four-manifold $M$, its quadratic two-type consists of $\pi=\pi_1(M)$, the deck-transformation module $A=\pi_2(M)$, its equivariant intersection pairing $\lambda_M:A\times A\to\mathbb Z[\pi]$, and the first Postnikov invariant $\kappa_M\in H^3(\pi;A)$. The latter class determines the fibration of the second Postnikov stage over $K(\pi,1)$ with fiber $K(A,2)$.

The equivariant pairing records all oriented intersections of lifted sphere representatives with translates in the universal cover, with coefficients indexed by $\pi$. It is Hermitian for the group-ring involution $g\mapsto g^{-1}$. An isomorphism of quadratic two-types consists of a group isomorphism and an equivariant module isomorphism carrying both this pairing and the Postnikov class to those of the other manifold.

Suppose $M$ and $N$ have finite fundamental groups and isomorphic quadratic two-types. Must $M$ and $N$ be homotopy equivalent? The problem concerns closed orientable manifolds, choosing orientations in order to compare the forms; it does not impose simple connectivity or assume that the finite groups are cyclic.

The desired homotopy equivalence is not required to induce a previously chosen isomorphism of the algebraic data. Nor does the conclusion claim homeomorphism or diffeomorphism. Replacing manifolds by arbitrary Poincaré complexes would change the realization constraints, and dropping finiteness of the fundamental group is outside the proposed implication. The question asks whether these data exhaust the homotopy information for the stated manifold class.
\subsection{Short English statement}
For orientable four-manifolds with finite fundamental group, does the quadratic two-type determine homotopy type?
\subsection{Sources}
\begin{itemize}
\item[S1] ULAM.AI and the UnsolvedMath Contributors, supplied problems.json, record 2931 (KP-4.55), Kirby's Problems in Low-Dimensional Topology. Catalogue identity and source transcription; CC BY 4.0.
\url{https://huggingface.co/datasets/ulamai/UnsolvedMath}.
\item[S2] K3: A New Problem List in Low-Dimensional Topology, Problem 4.55. Expanded formulation of the supplied catalogue statement.
\url{https://math.berkeley.edu/sites/default/files/surv-295-ruberman-watermarked-author-pdf.pdf}.
\end{itemize}
\end{document}
